\documentclass[a4paper]{article}
\usepackage{amsfonts}
\usepackage{amsmath}
\usepackage{amssymb}
\usepackage{graphicx}  

\newcommand{\degree}{^{\circ}}
\newcommand{\cis}{\operatorname{cis}}
\newcommand{\Bgsin}{\operatorname{Bgsin}}
\newcommand{\Bgcos}{\operatorname{Bgcos}}
\newcommand{\Bgtan}{\operatorname{Bgtan}}

\begin{document}
  
\noindent \large Auteur: Sadia Nazary  \\
\noindent \large Studierichting: Informatica\\
\noindent \large Oefeningenreeks 3C nr. 7\\

\medskip

\normalsize

Totale lengte prikkeldraad: $2400$ m \\

Noem $x$ de lengte van de zijde evenwijdig met het kanaal. \\

Voor de twee overige zijden blijft dan $2400-x$ meter draad over. \\

Elke van deze zijden heeft dus lengte $\dfrac{2400-x}{2}$. \\

De oppervlakte van de rechthoek is \\

$A(x)=x\cdot\dfrac{2400-x}{2}$ \\

$=1200x-\dfrac{x^2}{2}$ \\

$A'(x)=1200-x$ \\

$A'(x)=0$ \\

$1200-x=0$ \\

$x=1200$ \\

\medskip

\begin{center}
\begin{tabular}{c|ccc}
$x$ & $0$ & $1200$ & $2400$ \\
\hline
$A'(x)$ & $+$ & $0$ & $-$ \\
$A(x)$ & $\nearrow$ & & $\searrow$
\end{tabular}
\end{center}

$A''(x)=-1<0$ \\

zijde langs het kanaal $=1200$ m \\

andere twee zijden $=\dfrac{2400-1200}{2}=600$ m \\

\fbox{maximale oppervlakte $=720000\ \text{m}^2$}

\end{document}